% capitulos/cap04_geometria.tex
\chapter{La geometría de todas las barras}\label{cap:geometria}

En el \cref{cap:primer_notebook} calculamos la geometría de una barra, y en el
\cref{cap:describir_armadura} describimos la armadura completa con
diccionarios. Ahora unimos ambas cosas: calculamos la longitud y los cosenos
directores de \emph{todas} las barras, sin importar cuántas sean. La
herramienta que lo hace posible es el bucle, la instrucción que repite un
mismo cálculo para cada elemento de una colección. Con ella, nuestra
calculadora deja de depender del tamaño de la armadura.

\begin{objetivos}
  \item Repetir un cálculo con el bucle \py{for}.
  \item Recorrer un diccionario con \py{.items()} y acumular resultados.
  \item Calcular la geometría de todas las barras y guardarla en un
    diccionario nuevo, sin modificar los datos de entrada.
  \item Presentar los resultados en una tabla y dibujar cualquier armadura.
\end{objetivos}

% ---------------------------------------------------------------------------
\section{El mismo cálculo para cada barra}\label{sec:mismo_calculo}

La geometría de una barra depende solo de las coordenadas de sus dos nodos:
\[
  L = \sqrt{\Delta x^2 + \Delta y^2}, \qquad
  c = \frac{\Delta x}{L}, \qquad s = \frac{\Delta y}{L}.
\]
Para el ejemplo A, con tres barras, podríamos copiar tres veces el cálculo del
\cref{cap:primer_notebook}. Pero una armadura de cubierta típica tiene entre
20 y 50 barras, y la de un puente, varios cientos. Copiar el cálculo sería
tedioso y, sobre todo, una fuente segura de errores: basta con olvidar
cambiar un número en una de las copias.

Lo que necesitamos es escribir el cálculo \emph{una sola vez} y pedirle a
Python que lo repita para cada barra del diccionario \py{barras}.

% ---------------------------------------------------------------------------
\section{El bucle \texttt{for}}\label{sec:for}

Un \termino{bucle}{loop} \py{for} repite un bloque de instrucciones para cada
elemento de una colección. En cada repetición, la variable que sigue a
\py{for} toma el valor del elemento siguiente:
\fragmento{c04_for_lista}
\salida{c04_for_lista}

\begin{python}
El bloque que se repite es el que está \emph{sangrado} (desplazado hacia la
derecha) debajo de la línea del \py{for}, que termina en dos puntos. Python
usa la sangría, normalmente de cuatro espacios, para saber qué instrucciones
pertenecen al bucle. Colab la agrega automáticamente al pulsar Enter después
de los dos puntos; para salir del bloque basta con borrarla.
\end{python}

Si recorremos un diccionario, la variable toma, una por una, sus
\emph{claves}. Con cada clave podemos leer el valor correspondiente:
\fragmento{c04_for_dict}
\salida{c04_for_dict}

Casi siempre necesitamos la clave y el valor a la vez. El método
\py{.items()} entrega ambos en cada repetición, y los recibimos en dos
variables separadas por una coma:
\fragmento{c04_items}
\salida{c04_items}

Un uso muy frecuente de los bucles es \emph{acumular} un resultado. Empezamos
con un valor inicial y, en cada repetición, le sumamos algo. El operador
\py{+=} suma al valor que la variable ya tenía: \py{a += b} equivale a
\py{a = a + b}. Así obtenemos, por ejemplo, la longitud total de las barras:
\fragmento{c04_acumulador}
\salida{c04_acumulador}

% ---------------------------------------------------------------------------
\section{La geometría de todas las barras}\label{sec:geometria_todas}

Ya tenemos todas las piezas. Recorremos las barras; para cada una, buscamos
sus nodos en el diccionario \py{nodos}, calculamos las proyecciones, la
longitud y los cosenos directores, y guardamos el resultado en un diccionario
nuevo, \py{geometria}, con el número de barra como clave:
\fragmento{c04_geometria}

Este código forma parte de nuestra calculadora: es la primera pieza que
funciona para \emph{cualquier} armadura descrita con los diccionarios del
\cref{cap:describir_armadura}.

\begin{cuidado}
Guardamos los resultados en un diccionario nuevo, en lugar de agregarlos a
\py{barras}. Así los datos de entrada nunca cambian: si algo sale mal, podemos
volver a ejecutar el cálculo con la seguridad de partir de los mismos datos.
Mezclar datos y resultados en la misma estructura es una causa frecuente de
errores difíciles de rastrear.
\end{cuidado}

Para presentar los resultados recorremos \py{geometria} e imprimimos una fila
por barra. En las f-strings, el número antes del punto indica el
\emph{ancho} de la columna: \py{:7.3f} reserva siete caracteres con tres
decimales, y \py{:>5} alinea el número de barra a la derecha en cinco
caracteres. Así las columnas quedan alineadas:
\fragmento{c04_tabla}
\salida{c04_tabla}

% ---------------------------------------------------------------------------
\section{Verificación}\label{sec:verificacion_geometria}

Los valores coinciden con los que calculamos a mano en el
\cref{cap:primer_notebook}: la barra~1 es horizontal ($c = 1$, $s = 0$), y las
barras 2 y 3 son las hipotenusas de dos triángulos 3-4-5, con $L = 5$, $s =
0.8$ y $c = \pm 0.6$. El signo de $c$ en la barra~2 es negativo porque va del
nodo~2 al nodo~3, hacia la izquierda.

Una verificación que no depende del ejemplo es que, para toda barra, debe
cumplirse $c^2 + s^2 = 1$:
\fragmento{c04_verificacion}
\salida{c04_verificacion}

Finalmente, comprobamos que los datos de entrada no cambiaron: la barra~1
conserva sus cuatro datos originales y nada más.
\fragmento{c04_datos_intactos}
\salida{c04_datos_intactos}

% ---------------------------------------------------------------------------
\section{Dibujar cualquier armadura}\label{sec:dibujar_cualquier}

En el \cref{cap:describir_armadura} escribimos una línea \py{ax.plot} por
cada barra. Con un bucle, el mismo dibujo funciona para cualquier armadura.
Aprovechamos además para escribir el número de cada barra en su punto medio,
con \py{ax.text(x, y, texto)}:
\fragmento{c04_dibujo}
\grafico{c04_dibujo}

\begin{problemaresuelto}[La geometría del ejemplo B]
Calcular la longitud y los cosenos directores de las seis barras del ejemplo
B (\cref{fig:ejemplo_b}) y verificar los resultados.

\solucion
Aplicamos el mismo bucle a los diccionarios \py{nodos_b} y \py{barras_b} del
\cref{cap:describir_armadura}, y guardamos los resultados en
\py{geometria_b}:
\fragmento{c04_problema_resuelto}
\salida{c04_problema_resuelto}

\textbf{Verificación.} Las barras 1 y 3 son horizontales de \SI{3}{\metre} y
las barras 2 y 4, verticales de \SI{4}{\metre}. Las diagonales 5 y 6 miden
\SI{5}{\metre}; la diagonal~5 sube hacia la derecha ($c > 0$) y la 6, hacia la
izquierda ($c < 0$), como se ve en la figura.
\end{problemaresuelto}

Para resolver el problema copiamos el bucle y cambiamos los nombres de los
diccionarios. Funciona, pero es justamente la clase de repetición que
queríamos evitar. En el \cref{cap:rigidez_barra} aprenderemos a darle un
nombre a un cálculo, con una \emph{función}, para usarlo con cualquier
armadura sin copiarlo.

% ---------------------------------------------------------------------------
\repasoyresumen

\begin{itemize}
  \item El bucle \py{for} repite el bloque sangrado debajo de él para cada
    elemento de una colección.
  \item Al recorrer un diccionario, \py{for} entrega sus claves;
    \py{.items()} entrega la clave y el valor a la vez.
  \item Un acumulador empieza con un valor inicial y crece en cada repetición
    con \py{+=}.
  \item La geometría de todas las barras se guarda en un diccionario nuevo,
    \py{geometria}, con el número de barra como clave. Los datos de entrada no
    se modifican.
  \item En las f-strings, \py{:7.3f} fija el ancho y los decimales de una
    columna, y \py{:>5} alinea a la derecha.
  \item Para toda barra debe cumplirse $c^2 + s^2 = 1$.
\end{itemize}

% ---------------------------------------------------------------------------
\ejercicios

\begin{ejercicio}{4.1}
Modificar el bucle de la geometría para calcular también el ángulo $\theta$
de cada barra del ejemplo A, en grados, y mostrarlo para las tres barras.
\end{ejercicio}

\begin{ejercicio}{4.2}
El peso específico del acero es $\gamma = \SI{78.5}{\kilo\newton\per\metre\cubed}$.
Calcular el peso propio total de las barras del ejemplo A, sabiendo que el
peso de una barra es $\gamma A L$.
\end{ejercicio}

\begin{ejercicio}{4.3}
Calcular la longitud total de las barras del ejemplo B con un acumulador.
\end{ejercicio}

\begin{ejercicio}{4.4}
Dibujar la armadura del ejemplo B con un bucle, numerando sus barras, y
comparar el dibujo con la \cref{fig:ejemplo_b}.
\end{ejercicio}

\puntodecontrol{04}
